\chapter{Risultati e discussione}
\label{cap:risultati}

Questo capitolo raccoglie le considerazioni maturate nei vari capitoli di questo lavoro e rilegge le misure del Capitolo~\ref{cap:test} alla luce delle domande poste all'inizio, traendone due conclusioni. La prima riguarda ciò che è stato costruito: un sistema funzionante e misurato, che risponde a un bisogno concreto e potrebbe essere impiegato così com'è. La seconda, che costituisce il baricentro della tesi, riguarda l'AI generativa: un lavoro nato come esplorazione del suo apporto al rilevamento delle PII ne ridimensiona il ruolo sulla base dei dati raccolti durante le sperimentazioni.

% =====================================================================
\section{Il sistema realizzato}
\label{sec:risultati-sistema}
% =====================================================================

Al netto della questione generativa, il lavoro ha prodotto un \textbf{prototipo funzionante e misurato} che copre l'intera catena dal documento all'esito di conformità. A partire da una cartella di documenti non strutturati, il sistema ne estrae il testo (§\ref{sec:impl-b3-estrazione}), vi rileva le categorie di dati personali fondendo pattern e NER (§\ref{sec:impl-detection}), registra le PII trovate come \emph{riferimenti} minimizzati --- tipo, posizione, provenienza, mai il valore (§\ref{sub:dati-registro-pii}) ---, associa i documenti alle attività dichiarate nel registro dei trattamenti (§\ref{sec:b6-documento-trattamento}) e restituisce un esito di conformità leggibile dal DPO: categorie orfane, coperte, mancanti e superamento dei termini di conservazione (§\ref{sec:impl-b7-conformita}). È esattamente l'insieme degli obiettivi enunciati all'inizio (§\ref{sec:obiettivi}), realizzato end-to-end e operabile da un'interfaccia web unica (§\ref{sec:b8-dashboard}).

Sul metro sperimentale, la parte di rilevamento coglie quasi tutto ciò che deve trovare. La pipeline tradizionale raggiunge una \textbf{recall di 0.99} sul corpus aziendale, misurata sull'insieme completo dei $300$ documenti e delle $2918$ occorrenze annotate (Tabella~\ref{tab:test-pattern}, Tabella~\ref{tab:test-ner}). Il degrado misurato fra testo pulito e testo estratto dai file reali è nullo (Tabella~\ref{tab:test-estrazione}), sicché su documenti \emph{born-digital} lo stadio di estrazione si può considerare trasparente. La precisione è invece più bassa ($0.66$ sul campione su cui è tabulata la pipeline fusa, Tabella~\ref{tab:test-detection-confronto}), con una causa circoscritta a una sola categoria: il riconoscimento dei \emph{nomi di persona}, dove l'\textit{encoder} NER etichetta come tali molte espressioni che non lo sono, raccoglie da solo circa due terzi di tutti i falsi positivi, proporzione che si conferma sia sul campione sia sull'intero corpus. Il difetto è reale ma di natura benigna per l'uso previsto, poiché produce lavoro di revisione per il DPO anziché rischi di conformità non visti, ed è coerente con l'impostazione \textit{recall-oriented} del \emph{merge} (§\ref{sub:b4-merge}), che preferisce deliberatamente un falso positivo a una PII mancata. Trattandosi di numeri ottenuti su corpus sintetici, essi vanno letti come \emph{limite superiore} (§\ref{sub:test-estrazione}) più che come garanzia sul campo; bastano però a stabilire che l'impianto scelto --- Presidio con GLiNER, fusi da un motore di \textit{merge} \textit{recall-oriented} --- è adeguato al compito per cui è stato adottato, e a indicare con precisione dove andrebbe migliorato.

Oltre all'accuratezza, il sistema risponde a un \textbf{bisogno normativo concreto e mal coperto dal mercato}. L'obbligo di sapere dove stanno i dati personali, a quale trattamento si riferiscono e se sono ancora entro i termini di conservazione discende direttamente dalla nLPD (art.~6(4)) e dal GDPR (art.~5(1)(e), art.~30), mentre nella pratica gran parte di quei dati vive dispersa in documenti non strutturati, fuori da ogni visibilità sistematica (§\ref{sec:progetto-assegnato}). Gli strumenti commerciali esistenti sono tarati su GDPR o CCPA e non trattano nativamente le categorie della nLPD svizzera, e solo i prodotti \textit{enterprise} più avanzati integrano il registro dei trattamenti come riferimento dichiarativo. Il differenziale del sistema qui realizzato consiste proprio in quel collegamento: le PII rilevate sono \emph{confrontate con il dichiarato} e ne deriva un esito di conformità, che è il valore aggiunto per il DPO.

Il sistema è inoltre \textbf{impiegabile già oggi}, senza prerequisiti onerosi. Gira su hardware \textit{commodity} in un \textit{container} portabile (§\ref{sec:containerizzazione}), non richiede né GPU né addestramento su dati dell'organizzazione (§\ref{sec:considerazioni-training}), è coerente con la minimizzazione del dato per requisito --- i valori delle PII non vengono mai persistiti (§\ref{sub:dati-registro-pii}) --- e rileva a un costo contenuto in tempo ed energia (${\sim}11.8$~s e ${\sim}0.066$~Wh a documento, Tabella~\ref{tab:test-ai-benchmark}), il che lo rende eseguibile con continuità sull'intero patrimonio documentale. Restano naturalmente i limiti dichiarati --- assenza di OCR e di analisi di \textit{layout}, un solo profilo utente, corpus di valutazione sintetici (Capitolo~\ref{cap:sviluppi-futuri}) --- nessuno dei quali ostacola un impiego reale nello scenario per cui il sistema è stato pensato: documenti nativi digitali di un'organizzazione soggetta alla normativa sulla privacy.

% =====================================================================
\section{Costo e apporto dell'AI generativa}
\label{sec:risultati-ai}
% =====================================================================

Il secondo risultato riguarda l'apporto dell'AI generativa, e ne ridimensiona il ruolo rispetto a un'aspettativa largamente condivisa. Il modello generativo è oggi presentato, ben oltre l'ambito tecnico, come soluzione pressoché universale ai problemi di trattamento del testo, da preferire per principio agli approcci tradizionali. Applicata al rilevamento sistematico delle PII, tale aspettativa non trova riscontro nei dati raccolti: sul compito qui studiato l'impiego generativo comporta un costo superiore di un ordine di grandezza, in tempo e in energia, senza un corrispondente guadagno di qualità.

La conclusione arriva da due direzioni convergenti. Da un lato la letteratura la anticipa in termini generali (§\ref{sec:impatto-ambientale}): a parità di compito, un modello discriminativo consuma circa 0.002~kWh ogni 1000 inferenze contro le ${\sim}0.047$~kWh di un generativo, e sul piano economico il divario fra un \textit{encoder} e un modello generativo di frontiera arriva a rapporti dell'ordine di 300:1. Dall'altro --- ed è il contributo proprio di questo lavoro --- la misura diretta sul sistema reale conferma il divario e ne fissa l'entità sul \emph{nostro} hardware e sul \emph{nostro} compito. Il rilevamento tradizionale opera a ${\sim}11.8$~s e ${\sim}0.066$~Wh per documento; ogni passata generativa costa fra $40$ e $83$~s e fra $0.41$ e $0.76$~Wh per documento (Tabella~\ref{tab:test-ai-benchmark}): da ${\sim}3$ a ${\sim}7\times$ più lenta, da ${\sim}6$ a ${\sim}12\times$ più energivora. Va osservato che il divario, misurato su documenti reali, risulta inferiore a quello che la letteratura riporta su singole inferenze: entrambi i termini crescono con la lunghezza del testo, e su documenti estesi anche il rilevamento tradizionale sostiene il costo dell'\textit{encoder} NER sull'intero contenuto.

A tale costo non corrisponde un guadagno di qualità \emph{per il sistema}. La pipeline tradizionale raggiunge F1 $0.79$ sul corpus aziendale, e l'aggiunta del parere generativo la lascia invariata: le configurazioni di fusione --- le uniche operative, poiché il rilevatore generativo affianca i rilevatori tradizionali anziché sostituirli --- si collocano fra $0.76$ e $0.79$, ossia al più eguagliano la baseline (Tabella~\ref{tab:test-detection-confronto}). Il meccanismo è quello atteso: la fusione recupera occorrenze e porta la recall a $1.00$, introducendo però falsi positivi che riportano la precisione a $0.61$ e ne riassorbono il vantaggio.

Preso isolatamente, il rilevatore generativo restituisce un quadro meno netto, che va riportato con esattezza. Dei due modelli misurati uno ottiene una F1 superiore a quella della baseline, l'altro le resta sensibilmente al di sotto: da due osservazioni discordi si ricava che l'esito \textbf{dipende dal modello}, non una proprietà generale dell'approccio generativo, mentre il costo resta superiore di un ordine di grandezza per entrambi. Va inoltre osservato il modo in cui il primo ottiene il proprio vantaggio: la sua recall è $0.85$ contro lo $0.99$ della baseline, ossia rileva \emph{meno} PII, e la F1 gli risulta favorevole soltanto perché quell'indicatore pesa allo stesso modo precisione e recall. La pesatura non corrisponde alle priorità del dominio: una PII non rilevata è un rischio di conformità che resta invisibile, mentre un falso positivo si traduce in lavoro di revisione per il DPO. Il sistema riconosce tale asimmetria adottando un \emph{merge} a impostazione \textit{recall-oriented} (§\ref{sub:b4-merge}); sulla dimensione che per questo compito conta di più il rilevatore generativo resta perciò indietro.

Neppure la taglia del modello incide favorevolmente, e in questo caso il vincolo si manifesta prima ancora della qualità: il modello da 12B della stessa famiglia \textbf{non è risultato eseguibile} nel bersaglio di \textit{deployment} dichiarato. La memoria che esso occupa e quella dell'\textit{encoder} NER eccedono insieme quanto una macchina \textit{commodity} da 16\,GB può garantire, e tre configurazioni successive hanno perso fra il 27\% e il 40\% delle finestre di testo per interruzione del processo di inferenza (§\ref{sec:test-ai-multimodello}). La fascia superiore è perciò esclusa dal confronto quantitativo, con una precisazione che ne delimita la portata: non si afferma che un modello da 12B produrrebbe risultati peggiori, bensì che nel perimetro che il progetto si è dato --- esecuzione locale su hardware ordinario, imposta dalla minimizzazione del dato e non da una preferenza --- quel modello non è disponibile. Analoga esclusione vale per il modello \textit{reasoning}, che presenta la combinazione meno favorevole fra costo e risultato (§\ref{sec:test-ai-multimodello}).

È stata infine verificata l'ipotesi che a penalizzare l'AI sia la sola esecuzione su CPU, e il controllo l'ha esclusa: su GPU la medesima inferenza è ${\sim}9\times$ più veloce e di qualità identica (§\ref{sub:test-gpu}), poiché l'accelerazione riguarda la velocità del calcolo e lascia intatta la sua accuratezza. Quell'accelerazione non è d'altronde disponibile nel bersaglio di \textit{deployment} containerizzato, dove la GPU non è esposta alla macchina virtuale.

La \textbf{portata} di questa conclusione va delimitata con precisione. Non si tratta di un rifiuto dell'AI generativa in quanto tale, bensì della constatazione che, \emph{su questo compito e su questo tipo di dato}, il suo costo non risulta compensato. Il compito è il rilevamento \emph{sistematico e ricorrente} --- ogni documento, ripetutamente --- ed è la ricorrenza a moltiplicare il costo unitario fino a renderlo decisivo (§\ref{sec:impatto-ambientale}). La grandezza determinante è il \textbf{tempo} più che la spesa in valore assoluto, che resta di pochi centesimi: a decine di secondi per documento, una scansione che per via tradizionale richiede minuti si estende a giorni. Il dato in esame è inoltre regolare e nativo digitale, ossia il regime più favorevole ai rilevatori tradizionali, nel quale pattern e NER sono già prossimi al massimo delle loro prestazioni e lasciano all'AI un margine di scoperta minimo. Le condizioni nelle quali il parere generativo potrebbe invece risultare vantaggioso --- testo rumoroso o proveniente da OCR, formati inusuali, PII non strutturate che né i pattern né il NER prevedono --- sono esattamente quelle che i corpus impiegati non riproducono, e restano una questione aperta (§\ref{sec:sf-direzioni}). La scelta architetturale adottata è coerente con questa lettura: l'AI è mantenuta e collocata dove il suo costo non si moltiplica per il corpus --- uno slot \textbf{opzionale e campionato} nel rilevamento (§\ref{sec:impl-detection}) e i pochi compiti \emph{selettivi} a bassa frequenza (§\ref{sec:scelta-modello-llm}) --- riservandola ai casi in cui le interrogazioni sono poche e la comprensione semantica apporta un valore che gli altri approcci non possono dare.

La lezione metodologica va oltre il singolo sistema. Di fronte a un problema di trattamento del testo la domanda pertinente non è se l'AI generativa lo risolva, poiché quasi sempre la risposta è affermativa, ma se lo risolva \emph{meglio} degli approcci più semplici, in misura sufficiente a giustificarne il costo alla scala e alla frequenza d'uso previste. Per il rilevamento sistematico delle PII, misure alla mano, la risposta è negativa: la via tradizionale è sufficientemente accurata, più veloce di un ordine di grandezza e sostenibile con continuità.

% =====================================================================
\section{Considerazioni conclusive}
\label{sec:risultati-chiusura}
% =====================================================================

I due risultati si intersecano. Il sistema realizzato dimostra che la conformità alla nLPD sul contenuto documentale reale è affrontabile immediatamente, con mezzi leggeri e riproducibili; lo studio sull'AI generativa spiega \emph{perché} quei mezzi sono leggeri: la leggerezza è un esito misurato e non una rinuncia, poiché la via più costosa non avrebbe reso di più. Il contributo di questo lavoro consiste dunque in un prototipo funzionante e, insieme, in una risposta quantitativa e documentata a una domanda che spesso si dà per scontata: sul rilevamento sistematico dei dati personali l'approccio tradizionale è la scelta migliore.

Entrambi i risultati poggiano su un artefatto che il lettore può esaminare: il codice del sistema, i generatori dei corpus e la procedura che produce le misure discusse in questi due capitoli sono pubblicati in un repository ad accesso aperto \cite{algisi2026piidiscovery}, alla versione consegnata insieme a questa tesi. La scelta è coerente con il criterio che il Capitolo~\ref{cap:stato-arte} ha usato per giudicare gli strumenti esistenti (§\ref{sub:quadro-comparativo}): un lavoro che rimprovera alle suite commerciali di non essere ispezionabili deve, per primo, esporre il proprio funzionamento alla verifica.
